% Histogramme des volumes équivalents obtenus par 50 binômes (exercice 4, chapitre 27)
\begin{tikzpicture}
\begin{axis}[otMesures,otAxesBout,width=118mm,height=48mm,xmin=11.6,xmax=13.7,ymin=0,ymax=17,
  xtick={11.6,11.8,...,13.6},ytick={0,5,10,15},minor ytick={1,2,...,16},grid=none,ymajorgrids=true,yminorgrids=true,
  minor grid style={otRule!45,line width=0.2pt},
  xticklabel style={/pgf/number format/fixed,/pgf/number format/fixed zerofill,/pgf/number format/precision=1},
  xlabel={$V_\mathrm{E}$ (mL)},ylabel={effectif}]
\addplot[ybar interval,fill=otNavy!22,draw=otNavy,line width=0.5pt] coordinates {(11.7,2) (11.8,7) (11.9,15) (12.0,14) (12.1,8) (12.2,3) (12.3,0)};
\addplot[ybar interval,fill=otRed!25,draw=otRed!80!black,line width=0.5pt] coordinates {(13.4,1) (13.5,0)};
\node[pcLegende,anchor=south,text=otRed!80!black] at (axis cs:13.45,1.3) {1 binôme};
\end{axis}
\end{tikzpicture}
\par\smallskip\legende{Volumes équivalents obtenus par les 50 binômes d'un lycée pour le titrage d'une même solution, regroupés par classes de \qty{0,1}{mL}.}
